% Definitions and frozen diagnostics from v4.9.14, v4.9.15 and v4.9.16.
\subsection{A GP calculation of scan-global significance}
\label{v5-sec:global-method}\label{sec:v5-global}
A local probability describes a fluctuation at one fixed mass. A global
probability asks how often a complete background experiment would produce
an equally extreme result anywhere in the declared search. Nearby mass
hypotheses reuse counts and sidebands, so their fluctuations are correlated.
The significance-GP method of Ananiev and Read~\cite{AnanievRead2023}
estimates those correlations from the analysis response, then generates
many inexpensive correlated fields. It complements the look-elsewhere
approach described above and can retain nonstationary correlations and
changes in the contributing datasets.

The GP used for this purpose is distinct from the GP used to estimate the
mass background. Let $\psi$ denote the signal yield for an individual
dataset or the shared coupling parameter $\epsilon^2$ for a combination.
The auxiliary best fit $\widehat\psi$ is unconstrained in sign, so that
deficits remain visible; the physical discovery test still tests a
nonnegative signal. Let $r(n;m)$ denote the signed local significance mapping
from the profile likelihood at mass $m$:
\begin{equation}
 r(n;m)=\operatorname{sgn}(\widehat\psi)
 \sqrt{2[\ell(0,\widehat\theta_0)-\ell(\widehat\psi,\widehat\theta)]}.
 \label{v5-eq:signed-local-mapping}
\end{equation}
It expresses the likelihood improvement on a significance-like scale.
The conventional one-sided excess display uses $Z_{\rm local}=\max(0,r)$
and the asymptotic mapping $p_0=1-\Phi(Z_{\rm local})$; a negative value of
$r$ records a deficit rather than evidence for a positive signal.

For one coherent expected background spectrum $B$, the calculation scans
the unfluctuated spectrum and spectra with one support bin increased by
its Poisson standard deviation. It records
\begin{align}
 a_m&=r(B;m), & D_{im}&=r(B+\sqrt{B_i}e_i;m)-a_m,\\
 C_{mn}&=\sum_iD_{im}D_{in}, & s_m&=\sqrt{C_{mm}},
 & K_{mn}&=C_{mn}/(s_ms_n).
 \label{v5-eq:field-response}
\end{align}
A simulated significance field is then
$r_m^*=a_m+s_mz_m^*$ with $z^*\sim N(0,K)$. The nonzero offset $a_m$ and
nonunit width $s_m$ are explicit extensions of the ideally centered,
unit-width construction. Subtracting $a_m$ in the response calculation
centers the covariance estimate; it does not subtract events from data.

The declared principal ordering first maps each mass to its conditional
local probability,
\begin{equation}
 p_m(r)=\begin{cases}
 1,&r\leq0,\\
 1-\Phi[(r-a_m)/s_m],&r>0,
 \end{cases}\qquad
 T=\max_{m:r_m>0}\frac{r_m-a_m}{s_m}.
 \label{v5-eq:global-ordering}
\end{equation}
If no raw fitted signal is positive, $T=-\infty$. The global tail is the
fraction of complete simulated fields with $T^*\geq T_{\rm obs}$.
A separate comparison orders the scan by $\max_m\max(0,r_m)$.
These are different tests because the first accounts for a mass-dependent
background offset before ranking positive fluctuations. Their probabilities
must not be chosen after seeing which is more interesting. The empirical
positive-signal gate gives $p=1$ for nonpositive $r$, whereas the older
asymptotic display gives $p_0=0.5$; this convention explains some abrupt
changes between plots.

Each individual study began with ten complete Poisson scans, followed by
1,000 independently generated complete scans and 200,000 GP fields.
Every mass in one scan uses the same full spectrum. Sideband targets,
count-dependent errors and posterior predictions are updated; the reviewed
kernel coordinates remain fixed. Independently generated pointwise toy
collections cannot be joined by their row numbers to reconstruct this
correlation. The 2015, 2016 and 2021 grids contain 72, 142 and 201 integer
mass hypotheses over 19--90, 39--180 and 50--250~MeV, respectively.

The combined study evaluates the actual shared-coupling likelihood over
232 integer masses from 19 to 250~MeV. The active datasets are 2015 at
19--38~MeV, 2015+2016 at 39--49~MeV, all three at 50--90~MeV,
2016+2021 at 91--180~MeV, and 2021 at 181--250~MeV. Independent
constituent random streams are paired into 1,000 joint experiments, with
each whole experiment preserved across the grid. The response basis uses
1,626 support bins and 1,627 deterministic scans. Drawing one joint field
preserves correlations across membership boundaries; combining separate
p-values or treating segments as independent would define another test.

\begin{table}[htbp]\centering\small
\begin{tabular}{lrrrr}\toprule
Profiled scan & Selected mass & GP exceedances & GP global $p$ & Direct exceedances\\
 & [MeV] & out of 200,000 & & out of 1,000\\\midrule
2015 & 51 & 3,797 & 0.018985 & 13\\
2016 & 42 & 0 & $<1.498\times10^{-5}$ & 0\\
2021 10\% & 92 & 5,038 & 0.02519 & 26\\
Combined union & 76 & 0 & $<1.498\times10^{-5}$ & 0\\\bottomrule
\end{tabular}
\caption{Existing conditional significance-GP diagnostics under the
archived coherent stress background, using the principal ordering in
Eq.~\ref{v5-eq:global-ordering}. Zero GP exceedances give the displayed
one-sided 95\% sampling bound; zero direct exceedances give only
$p<0.002991$. These bounds condition on their respective simulation models.
The table reports neither a validated rare particle tail nor a global
calibration of the two-truth envelope in Appendix~\ref{v5-sec:calibration}.}
\label{v5-tab:global-diagnostics}
\end{table}

The latest probability audit gives a particularly instructive combined
example. At 76~MeV the observed $r=0.166$ corresponds to an asymptotic local
$p_0=0.434$. The stress background predicts $a_m=-8.700$ with width 0.979,
so its offset-adjusted score is 9.053. This large score measures a
mismatch with that selected background reference. It cannot be reported
as a $9\sigma$ particle signal. Under the alternative raw ordering, every
simulated scan exceeds the observed maximum and the estimated global
probability is one. Neither comparison is a general goodness-of-fit test.

Direct scans check central distributions, widths, correlations and
accessible maximum tails. In the 2016 profiled study one unadjusted
bulk-distribution comparison has nominal $p=0.034$, reinforcing the need
to inspect approximation discrepancies. Agreement elsewhere and the
absence of adjusted marginal flags do not validate the unresolved rare
tails. The 2016 stress shape also retains source-fit and numerical
qualifications discussed in Appendix~\ref{v5-sec:calibration}.

\paragraph{Result still to be supplied for final interpretation.}
Final particle-global results require a justified background family and
kernel policy, independent validation of local probabilities, and adequate
simulation of the complete chosen mass and model search. Until those
conditions are met, Table~\ref{v5-tab:global-diagnostics} retains the
conditional diagnostics and leaves unresolved tails as bounds.
